AI agents have significant potential to reshape cybersecurity, making a thorough assessment of their capabilities critical.
However, existing evaluations fall short, because they are based on small-scale benchmarks and only measure static outcomes, failing to capture the full, dynamic range of real-world security challenges.
To address these limitations, we introduce \sysname, a large-scale benchmark featuring 1,507 real-world vulnerabilities across 188 software projects.
Adjustable to different vulnerability analysis settings, \sysname primarily tasks agents with generating a proof-of-concept test that reproduces a vulnerability, given only its text description and the corresponding codebase.
Our extensive evaluation highlights that \sysname effectively differentiates agents' and models' cybersecurity capabilities.
Even the top-performing combinations only achieve a $\sim$20\% success rate, demonstrating the overall difficulty of \sysname. 
Beyond static benchmarking, we show that \sysname leads to the discovery of 34 zero-day vulnerabilities and 18 historically incomplete patches.
These results underscore that \sysname is not only a robust benchmark for measuring AI's progress in cybersecurity but also a platform for creating direct, real-world security impact.